\section{Where Budget Portability Fails}
\label{sec:portability}

\subsection{RQ1: Do Responses Transfer When Only the Graph Changes?}
Retrospective first-passing actions frequently fail after a graph-only rebuild, even though collection membership, queries, and truth are unchanged. Table~\ref{tab:graph-only} reports absolute target risk, the unresolved-reference risk, and their difference. Across hnswlib and Faiss on both datasets, the increment ranges from 17.17 to 23.60 percentage points; every query-cluster interval is above zero. This is a response-level result independent of TCP's pooling or recalibration.

\begin{table*}[t]
\caption{Graph-only response transfer at $k=10$, $\tau=0.95$. Each row has 24 builds, 552 directed pairs, and 750 shared query IDs. Risk entries and increments are percentages or percentage points [95\% query-cluster interval], conditional on the fixed build set. The reference uses the target's own first-passing response; its failures are unresolved queries. Finite variation counts queries with at least two distinct finite first-passing actions, divided by all 750 queries.}
\label{tab:graph-only}
\small\begin{tabular*}{\textwidth}{@{\extracolsep{\fill}}llcccr@{}}\toprule
Implementation & Dataset & Absolute risk [95\% CI] & Reference risk [95\% CI] & Increment [95\% CI] & Finite variation\\\midrule
\GraphOnlyRows
\bottomrule\end{tabular*}
\end{table*}

The variation is principally among finite budgets: 75.87--91.07\% of queries have different finite first-passing actions across builds, while only 0.53--2.67\% switch between resolved and unresolved states. Removing the 1\% of queries with largest incremental contribution leaves increments of 21.37/16.95 percentage points for hnswlib and 23.43/17.53 for Faiss, on SIFT/Arxiv respectively. Endpoint censoring and a handful of extreme queries therefore do not explain the result.

This experiment diagnoses the budget--response relation using retrospective truth on both source and target. It establishes that source success alone is insufficient evidence for action reuse. The measured mechanism is behavioral variation across builds; structural attribution to particular edges or navigation paths lies outside this response-level design.

\subsection{RQ2: Does an Executable Policy Deteriorate after Refresh?}
The matched old/new replay holds the nine-history policy fixed and applies it to both snapshots. Unlike RQ1, this changes 5\% of the vectors as well as rebuilding. SIFT's empirical risk rises from 3.77\% to 6.69\%; Arxiv-Nomic's rises from 3.19\% to 7.24\%. The corresponding paired increments [crossed 95\% interval] are 2.92 [1.75, 4.07] and 4.05 [2.86, 5.31] percentage points. Thus this is a measured deterioration, not merely an absolute target risk with an unknown source baseline.

Independent single-policy CP checks at $\alpha=0.05$ qualify six of ten old SIFT policies and five of ten old Arxiv policies. All eleven fail the corresponding target check. This paired diagnostic uses per-policy checks rather than simultaneous control over the eleven transitions. Target evaluation risk exceeds 5\% in all ten builds of each dataset; removing the highest-risk target leaves 6.61\% and 7.13\%. The deterioration is distributed across targets.

The same policy information is held fixed in each old/new comparison, but the experiment does not separate changed graph paths from changed nearest-neighbor membership. Together, RQ1 and RQ2 establish complementary facts: unchanged data do not make the response function portable, and a recurring-query policy can lose its old qualification after refresh plus rebuilding.
